\documentclass[11pt]{article}
\usepackage[margin=1in]{geometry}
\usepackage{amsmath,amssymb,amsthm}
\usepackage{dsfont}
\usepackage{xcolor}
\usepackage{booktabs}
\usepackage{array}
\usepackage{placeins}
\usepackage{hyperref}

\hypersetup{colorlinks=true, linkcolor=blue!60!black, citecolor=blue!60!black, urlcolor=blue!60!black}

\title{CMR NK-DSGE with Financial Frictions\\[4pt]
\large Corrected Model Documentation}
\author{Anna Smirnova}
\date{February 2026}

\begin{document}
\maketitle

\noindent This document is the definitive reference for the disaster model as
implemented in \texttt{deqn-jax}. It incorporates all errata corrections
(see \texttt{docs/disaster\_model\_errata.md}), all three analytical
eliminations, and the split Phillips curve formulation. When this document and
\texttt{disaster.tex} disagree, \emph{this document is correct}.

\bigskip\noindent
\textbf{Color code:}
\textcolor{green!60!black}{green $=$ state variables},
\textcolor{orange!80!black}{orange $=$ current-period controls},
\textcolor{blue}{blue $=$ next-period controls},
\textcolor{brown}{brown $=$ definitions},
\textcolor{magenta}{magenta $=$ next-period exogenous}.


%==========================================================================
\section{Variables}
%==========================================================================

\subsection{State variables (13)}

\begin{itemize}
\item \textbf{Endogenous lags (8):}\;
  $\textcolor{green!60!black}{\pi_{t-1}}$,\;
  $\textcolor{green!60!black}{k_{t-1}}$,\;
  $\textcolor{green!60!black}{c_{t-1}}$,\;
  $\textcolor{green!60!black}{q_{t-1}}$,\;
  $\textcolor{green!60!black}{i_{t-1}}$,\;
  $\textcolor{green!60!black}{R_{t-1}}$,\;
  $\textcolor{green!60!black}{\tilde{w}_{t-1}}$,\;
  $\textcolor{green!60!black}{L_{t-1}}$

\item \textbf{Exogenous (5):}\;
  $\textcolor{green!60!black}{\epsilon_t}$,\;
  $\textcolor{green!60!black}{\mu_{\Upsilon,t}}$,\;
  $\textcolor{green!60!black}{g_t}$,\;
  $\textcolor{green!60!black}{\mu_{z,t}}$,\;
  $\textcolor{green!60!black}{m_t^p}$
\end{itemize}

\subsection{Policy variables (11 network outputs)}

After three analytical eliminations ($s_t, L_t, \bar\omega_t$), the neural
network outputs:
\[
  \textcolor{orange!80!black}{\lambda_{z,t}},\quad
  \textcolor{orange!80!black}{i_t},\quad
  \textcolor{orange!80!black}{\pi_t},\quad
  \textcolor{orange!80!black}{\tilde{w}_t},\quad
  \textcolor{orange!80!black}{h_t},\quad
  \textcolor{orange!80!black}{F_{w,t}},\quad
  \textcolor{orange!80!black}{F_{p,t}},\quad
  \textcolor{orange!80!black}{q_t},\quad
  \textcolor{orange!80!black}{c_t},\quad
  \textcolor{orange!80!black}{K_{p,t}},\quad
  \textcolor{orange!80!black}{K_{w,t}}
\]

\noindent\textbf{Why $c_t$ is a network output.}\;
The resource constraint uniquely determines $c_t$, but eliminating it
analytically creates a singularity: when the resource constraint yields
$c_t \approx 0$, the habit term $c_t\mu_{z,t} - b\,c_{t-1} \to 0$ and
the consumption Euler diverges. With $c_t$ as a network output, the optimizer
directly controls consumption. The resource constraint (Eq.~9) enforces
consistency.

\medskip\noindent\textbf{Why $K_{p,t}$ and $K_{w,t}$ are network outputs.}\;
The Phillips curve composites $K_p$, $K_w$ were originally computed as
nonlinear functions of $(F_p, \pi, \tilde\pi, \ldots)$ involving extreme
exponents (${}^{-5}$, ${}^{-12}$). These exponent chains create near-perfect
gradient cancellation: the derivatives of $K_p$ and the recursion RHS with
respect to $\pi$ have nearly opposite signs and magnitudes, leaving
$|\partial r/\partial \pi| \approx 0$. This froze the price/wage ratios
during training.

Making $K_p$, $K_w$ network outputs and splitting each Phillips K equation
into a \emph{definition} (log-space) and a \emph{recursion} (level-space)
breaks this cancellation. The definition equation constrains $K_p$ to match
its analytical value; the recursion equation has clean gradient $\partial r/\partial K_p = 1$.

\subsection{Policy bounds}

\begin{table}[h]
\centering
\begin{tabular}{lcccl}
\toprule
Variable & Lower & Upper & SS value & Bounding \\
\midrule
$\lambda_{z,t}$ & 0.20 & $\infty$ & 0.6018 & softplus \\
$i_t$            & 0.40 & $\infty$ & 0.7947 & softplus \\
$\pi_t$          & 0.95 & 1.10     & 1.0123 & sigmoid \\
$\tilde{w}_t$    & 1.00 & $\infty$ & 1.9202 & softplus \\
$h_t$            & 0.60 & $\infty$ & 0.9440 & softplus \\
$F_{w,t}$        & 0.30 & $\infty$ & 0.8853 & softplus \\
$F_{p,t}$        & 2.00 & $\infty$ & 4.7359 & softplus \\
$q_t$            & 0.50 & $\infty$ & 1.0000 & softplus \\
$c_t$            & 0.50 & $\infty$ & 1.5937 & softplus \\
$K_{p,t}$        & 1.00 & $\infty$ & 4.8308 & softplus \\
$K_{w,t}$        & 0.50 & $\infty$ & 2.2071 & softplus \\
\bottomrule
\end{tabular}
\caption{Per-variable output bounding. Softplus:
$\text{output} = \ell + \log(1+e^{\text{raw}})$; no gradient death at the bound.
Sigmoid: $\text{output} = \ell + (u-\ell)\,\sigma(\text{raw})$; used only for
$\pi_t$ to prevent Phillips curve blowup from the ${}^{-5}$ exponent.}
\end{table}

\subsection{Definition variables (computed, not learned)}

\textcolor{brown}{$s_t$} (marginal cost),
\textcolor{brown}{$L_t$} (leverage),
\textcolor{brown}{$\bar\omega_t$} (default threshold),
\textcolor{brown}{$k_t$},
\textcolor{brown}{$n_t$},
\textcolor{brown}{$y_{z,t}$},
\textcolor{brown}{$c_t^{\text{implied}}$} (from resource constraint),
\textcolor{brown}{$y_t^{\text{GDP}}$},
\textcolor{brown}{$r_t^k$},
\textcolor{brown}{$R_t^k$},
\textcolor{brown}{$R_t$},
\textcolor{brown}{$\tilde\pi_t$},
\textcolor{brown}{$\tilde\pi_{w,t}$},
\textcolor{brown}{$\pi_{w,t}$},
\textcolor{brown}{$S(\cdot)$},
\textcolor{brown}{$S'(\cdot)$},
\textcolor{brown}{$F(\bar\omega_t)$},
\textcolor{brown}{$G(\bar\omega_t)$},
\textcolor{brown}{$\Gamma(\bar\omega_t)$},
\textcolor{brown}{$K_{p,\text{inner}}$},
\textcolor{brown}{$K_{w,\text{inner}}$}.

\medskip\noindent\textbf{Note:} $K_{p,t}$ and $K_{w,t}$ are network outputs,
not definitions. The ``inner'' terms $K_{p,\text{inner}}$ and
$K_{w,\text{inner}}$ are computed analytically and used only in the
definition equations (Eq.~2a, 4a).


%==========================================================================
\section{Definitions}
%==========================================================================

\paragraph{Inflation indexation.}
\begin{align*}
  \textcolor{brown}{\tilde\pi_t}
    &= \pi_{\text{ss}}^{\,\iota}\;\textcolor{green!60!black}{\pi_{t-1}}^{1-\iota}
  \\
  \textcolor{brown}{\tilde\pi_{w,t}}
    &= \pi_{\text{ss}}^{\,\iota_w}\;\textcolor{green!60!black}{\pi_{t-1}}^{1-\iota_w}
  \\
  \textcolor{brown}{\pi_{w,t}}
    &= \textcolor{orange!80!black}{\pi_t}\;
       \frac{\textcolor{orange!80!black}{\tilde{w}_t}}{\textcolor{green!60!black}{\tilde{w}_{t-1}}}
\end{align*}

\paragraph{Investment adjustment costs.}
\begin{align*}
  \textcolor{brown}{S\!\left(\tfrac{\mu_{z,t}\,i_t}{i_{t-1}}\right)}
    &= \frac{\kappa}{2}\!\left(\textcolor{green!60!black}{\mu_{z,t}}\,
       \frac{\textcolor{orange!80!black}{i_t}}{\textcolor{green!60!black}{i_{t-1}}}
       - \mu_{z,\text{ss}}\right)^{\!2}
  \\
  \textcolor{brown}{S'\!\left(\tfrac{\mu_{z,t}\,i_t}{i_{t-1}}\right)}
    &= \kappa\!\left(\textcolor{green!60!black}{\mu_{z,t}}\,
       \frac{\textcolor{orange!80!black}{i_t}}{\textcolor{green!60!black}{i_{t-1}}}
       - \mu_{z,\text{ss}}\right)
\end{align*}

\paragraph{Capital and factor prices.}
\begin{align*}
  \textcolor{brown}{k_t}
    &= (1-\delta)\,\frac{\textcolor{green!60!black}{k_{t-1}}}{\textcolor{green!60!black}{\mu_{z,t}}}
       + \bigl(1-\textcolor{brown}{S(\cdot)}\bigr)\,\textcolor{orange!80!black}{i_t}
  \\
  \textcolor{brown}{s_t}
    &= \frac{1}{\textcolor{green!60!black}{\epsilon_t}}
       \left(\frac{\textcolor{green!60!black}{\mu_{z,t}}\,\textcolor{orange!80!black}{h_t}}
       {\textcolor{green!60!black}{k_{t-1}}}\right)^{\!\alpha}
       \frac{\textcolor{orange!80!black}{\tilde{w}_t}}{1-\alpha}
  & &\text{(eliminates eq.~10)}
  \\
  \textcolor{brown}{r_t^k}
    &= \textcolor{green!60!black}{\epsilon_t}\,\alpha
       \left(\frac{\textcolor{green!60!black}{\mu_{z,t}}\,\textcolor{orange!80!black}{h_t}}
       {\textcolor{green!60!black}{k_{t-1}}}\right)^{\!1-\alpha}\textcolor{brown}{s_t}
  \\
  \textcolor{brown}{R_t^k}
    &= \frac{(1-\tau_k)\,\textcolor{brown}{r_t^k} + (1-\delta)\,\textcolor{orange!80!black}{q_t}}
       {\textcolor{green!60!black}{q_{t-1}}}\;\textcolor{orange!80!black}{\pi_t}
       + \tau_k\,\delta
\end{align*}

\paragraph{Financial frictions: lognormal distribution.}
\begin{align*}
  F(\bar\omega) &= \Phi\!\left(\frac{\ln\bar\omega + \tfrac{1}{2}\sigma_\omega^2}{\sigma_\omega}\right)
  &
  G(\bar\omega) &= \Phi\!\left(\frac{\ln\bar\omega - \tfrac{1}{2}\sigma_\omega^2}{\sigma_\omega}\right)
  \\
  \Gamma(\bar\omega) &= \bar\omega\,(1-F(\bar\omega)) + G(\bar\omega)
  &
  G'(\bar\omega) &= \frac{\phi\!\left(\frac{\ln\bar\omega - \tfrac{1}{2}\sigma_\omega^2}{\sigma_\omega}\right)}
    {\sigma_\omega\,\bar\omega}
  \\
  \Gamma'(\bar\omega) &= 1 - F(\bar\omega)
\end{align*}
where $\Phi$ is the standard normal CDF and $\phi$ is the standard normal PDF.

\paragraph{Default threshold $\bar\omega_t$ (eliminates eq.~8, Newton solver).}
Define the bank's expected return function:
\begin{equation*}
  h(\bar\omega) = \bar\omega\,(1-F(\bar\omega)) + (1-\mu)\,G(\bar\omega)
  = \Gamma(\bar\omega) - \mu\,G(\bar\omega)
\end{equation*}
and the target from the zero-profit condition:
\begin{equation*}
  \text{target}_t = \frac{\textcolor{green!60!black}{L_{t-1}} - 1}
    {\textcolor{green!60!black}{L_{t-1}}\;\textcolor{brown}{R_t^k}\,/\,\textcolor{green!60!black}{R_{t-1}}}
\end{equation*}
Solve $h(\bar\omega_t) = \text{target}_t$ by Newton iteration:
\begin{equation*}
  \boxed{\bar\omega^{(n+1)} = \bar\omega^{(n)}
    - \frac{h(\bar\omega^{(n)}) - \text{target}}{h'(\bar\omega^{(n)})},
  \qquad
  h'(\bar\omega) = 1-F(\bar\omega) - \mu\,G'(\bar\omega)}
\end{equation*}
Converges in ${\sim}2$ iterations from $\bar\omega_0 = 0.488$;
implementation uses 10 fixed iterations, $\bar\omega \in [0.01,\,3.0]$, via
\texttt{jax.lax.scan}.

\paragraph{Entrepreneur net worth.}
\begin{equation*}
  \textcolor{brown}{n_t} = \frac{\gamma_e}{\textcolor{orange!80!black}{\pi_t}\,\textcolor{green!60!black}{\mu_{z,t}}}
    \bigl(1 - \Gamma(\textcolor{brown}{\bar\omega_t})\bigr)\,
    \textcolor{brown}{R_t^k}\,\textcolor{green!60!black}{q_{t-1}}\,\textcolor{green!60!black}{k_{t-1}}
    + w^e
\end{equation*}
\textbf{Note:} The original code had $(1 - \mu\,G)$ instead of $(1-\Gamma)$;
see errata Bug~1.

\paragraph{Leverage (eliminates eq.~12).}
\begin{equation*}
  \textcolor{brown}{L_t} = \frac{\textcolor{orange!80!black}{q_t}\,\textcolor{brown}{k_t}}
    {\textcolor{brown}{n_t}}
\end{equation*}

\paragraph{Output and GDP.}
\begin{align*}
  \textcolor{brown}{y_{z,t}}
    &= \textcolor{green!60!black}{\epsilon_t}
       \left(\frac{\textcolor{green!60!black}{k_{t-1}}}{\textcolor{green!60!black}{\mu_{z,t}}}\right)^{\!\alpha}
       \textcolor{orange!80!black}{h_t}^{1-\alpha} - \Phi
  \\
  \textcolor{brown}{y_t^{\text{GDP}}}
    &= \textcolor{green!60!black}{g_t} + \textcolor{orange!80!black}{c_t}
       + \frac{\textcolor{orange!80!black}{i_t}}{\textcolor{green!60!black}{\mu_{\Upsilon,t}}}
\end{align*}
\textbf{Note:} $y_t^{\text{GDP}}$ uses the \emph{network output} $c_t$
(orange), not the implied consumption from the resource constraint.

\paragraph{Implied consumption (for Eq.~9 residual only).}
\begin{equation*}
  \textcolor{brown}{c_t^{\text{implied}}} = \textcolor{brown}{y_{z,t}} - \textcolor{green!60!black}{g_t}
    - \frac{\textcolor{orange!80!black}{i_t}}{\textcolor{green!60!black}{\mu_{\Upsilon,t}}}
    - \underbrace{\Theta\frac{1-\gamma_e}{\gamma_e}(\textcolor{brown}{n_t} - w^e)}_{\text{entrepreneur consumption}}
    - \underbrace{\frac{\mu\,G(\textcolor{brown}{\bar\omega_t})\,\textcolor{brown}{R_t^k}\,
      \textcolor{green!60!black}{q_{t-1}}\,\textcolor{green!60!black}{k_{t-1}}}
      {\textcolor{green!60!black}{\mu_{z,t}}\,\textcolor{orange!80!black}{\pi_t}}}_{\text{monitoring cost}}
\end{equation*}
This is \emph{not} used in any equation except Eq.~9 (resource constraint).
The consumption Euler (Eq.~5) uses the network output $\textcolor{orange!80!black}{c_t}$ directly.

\paragraph{Interest rate (Taylor rule).}
\begin{equation*}
  \textcolor{brown}{R_t}
    = R_{\text{ss}} \left(\frac{\textcolor{green!60!black}{R_{t-1}}}{R_{\text{ss}}}\right)^{\!\rho_p}
      \!\left[
        \left(\frac{\textcolor{orange!80!black}{\pi_t}}{\pi_{\text{ss}}}\right)^{\!\alpha_\pi}
        \left(\frac{\textcolor{brown}{y_t^{\text{GDP}}}}{y_{\text{ss}}}\right)^{\!\alpha_y}
      \right]^{1-\rho_p}
      \!\exp(\textcolor{green!60!black}{m_t^p})
\end{equation*}

\paragraph{Phillips curve inner terms.}
With \emph{soft floor}: $\lfloor x \rfloor_\varepsilon = \varepsilon + \frac{1}{s}\log(1+e^{s(x-\varepsilon)})$,\;
$\varepsilon = 0.01$, $s=10$.
\begin{align*}
  \textcolor{brown}{K_{p,\text{inner}}} &= \left\lfloor
    \frac{1 - \xi_p\bigl(\tilde\pi_t/\pi_t\bigr)^{1/(1-\lambda_f)}}{1-\xi_p}
  \right\rfloor_{0.01}
  \\
  \textcolor{brown}{K_{w,\text{inner}}} &= \left\lfloor
    \frac{1 - \xi_w\bigl(\tilde\pi_{w,t}\mu_{z,\text{ss}}/\pi_{w,t}\bigr)^{1/(1-\lambda_w)}}{1-\xi_w}
  \right\rfloor_{0.01}
\end{align*}
These are used only in the definition equations (Eq.~2a, 4a).
The actual $K_{p,t}$ and $K_{w,t}$ are network outputs.

\paragraph{Computation order.}
\begin{equation*}
  \underbrace{(\tilde\pi,\tilde\pi_w,\pi_w,S,S')}_{\text{from states + policies}}
  \;\to\; (k, s, r^k, R^k)
  \;\to\; \bar\omega_t\;\text{(Newton)}
  \;\to\; (F, G, \Gamma, n, L)
  \;\to\; (y_z, c^{\text{implied}}, y^{\text{GDP}}, R, K_{p,\text{inner}}, K_{w,\text{inner}})
\end{equation*}


%==========================================================================
\section{Equilibrium Equations (11)}
\label{sec:equations}
%==========================================================================

The 11 residual equations solved by the neural network.
The Phillips curve $K$ equations are each split into a \emph{definition}
(log-space, linearizes extreme exponents) and a \emph{recursion}
(level-space, clean unit gradient on $K$).

\setcounter{equation}{0}

\paragraph{Eq.~1: Price Phillips curve --- numerator ($F_p$).}
\begin{equation}
  \textcolor{orange!80!black}{\lambda_{z,t}}\,\textcolor{brown}{y_{z,t}}
  + \beta\,\xi_p\;\mathds{E}_t\!\left[
    \left(\frac{\textcolor{brown}{\tilde\pi_{t+1}}}{\textcolor{blue}{\pi_{t+1}}}\right)^{\!\frac{1}{1-\lambda_f}}
    \textcolor{blue}{F_{p,t+1}}
  \right]
  - \textcolor{orange!80!black}{F_{p,t}} = 0
\end{equation}

\paragraph{Eq.~2a: Price Phillips $K_p$ --- definition (log-space).}
Constrains the network's $K_{p,t}$ to match its analytical value:
\begin{equation}
  \log\textcolor{orange!80!black}{K_{p,t}}
  - \log\textcolor{orange!80!black}{F_{p,t}}
  - (1 - \lambda_f)\,\log\textcolor{brown}{K_{p,\text{inner}}} = 0
  \tag{2a}
\end{equation}

\paragraph{Eq.~2b: Price Phillips $K_p$ --- recursion (level-space).}
Gradient w.r.t.\ $K_{p,t}$ is exactly 1:
\begin{equation}
  \textcolor{orange!80!black}{K_{p,t}}
  - \textcolor{orange!80!black}{\lambda_{z,t}}\,\lambda_f\,
    \textcolor{brown}{y_{z,t}}\,\textcolor{brown}{s_t}
  - \beta\,\xi_p\;\mathds{E}_t\!\left[
    \left(\frac{\textcolor{brown}{\tilde\pi_{t+1}}}{\textcolor{blue}{\pi_{t+1}}}\right)^{\!\frac{\lambda_f}{1-\lambda_f}}
    \textcolor{blue}{K_{p,t+1}}
  \right] = 0
  \tag{2b}
\end{equation}

\paragraph{Eq.~3: Wage Phillips curve --- numerator ($F_w$).}
\begin{equation}
  \frac{\textcolor{orange!80!black}{h_t}(1-\tau_l)}{\lambda_w}\,\textcolor{orange!80!black}{\lambda_{z,t}}
  + \beta\,\xi_w\;\mathds{E}_t\!\left[
    \textcolor{magenta}{\mu_{z,t+1}}^{\frac{\iota_\mu}{1-\lambda_w}-1}
    \mu_{z,\text{ss}}^{\frac{1-\iota_\mu}{1-\lambda_w}}
    \frac{(\textcolor{brown}{\tilde\pi_{w,t+1}})^{\frac{1}{1-\lambda_w}}}
         {(\textcolor{brown}{\pi_{w,t+1}})^{\frac{\lambda_w}{1-\lambda_w}}}
    \frac{\textcolor{blue}{F_{w,t+1}}}{\textcolor{blue}{\pi_{t+1}}}
  \right]
  - \textcolor{orange!80!black}{F_{w,t}} = 0
\end{equation}

\paragraph{Eq.~4a: Wage Phillips $K_w$ --- definition (log-space).}
\begin{equation}
  \log\textcolor{orange!80!black}{K_{w,t}}
  - \log\!\left(\frac{1}{\psi_L}\right)
  - \bigl(1 - \lambda_w(1+\sigma_L)\bigr)\,\log\textcolor{brown}{K_{w,\text{inner}}}
  - \log\textcolor{orange!80!black}{\tilde{w}_t}
  - \log\textcolor{orange!80!black}{F_{w,t}} = 0
  \tag{4a}
\end{equation}

\paragraph{Eq.~4b: Wage Phillips $K_w$ --- recursion (level-space).}
\begin{equation}
  \textcolor{orange!80!black}{K_{w,t}}
  - \textcolor{orange!80!black}{h_t}^{1+\sigma_L}
  - \beta\,\xi_w\;\mathds{E}_t\!\left[
    \left(\frac{\textcolor{brown}{\tilde\pi_{w,t+1}}\,\mu_{z,\text{ss}}}
         {\textcolor{brown}{\pi_{w,t+1}}}\right)^{\!\frac{\lambda_w}{1-\lambda_w}(1+\sigma_L)}
    \!\textcolor{blue}{K_{w,t+1}}
  \right] = 0
  \tag{4b}
\end{equation}

\paragraph{Eq.~5: Consumption Euler (multiply-through formulation).}
The standard formulation has $\mu_z / \text{habit}$ singularities.
Multiplying through by both habit terms eliminates all divisions:
\begin{equation}
  (1+\tau_c)\,\textcolor{orange!80!black}{\lambda_{z,t}}\,
  \textcolor{brown}{\text{habit}_{t}}\,\textcolor{brown}{\text{habit}_{t+1}}
  - \textcolor{green!60!black}{\mu_{z,t}}\,\textcolor{brown}{\text{habit}_{t+1}}
  + \beta\,b\,\textcolor{brown}{\text{habit}_{t}} = 0
\end{equation}
where
\begin{align*}
  \textcolor{brown}{\text{habit}_t}
    &= \bigl\lfloor \textcolor{orange!80!black}{c_t}\,\textcolor{green!60!black}{\mu_{z,t}}
       - b\,\textcolor{green!60!black}{c_{t-1}} \bigr\rfloor_{0.01}
  \\
  \textcolor{brown}{\text{habit}_{t+1}}
    &= \bigl\lfloor \textcolor{blue}{c_{t+1}}\,\textcolor{magenta}{\mu_{z,t+1}}
       - b\,\textcolor{orange!80!black}{c_t} \bigr\rfloor_{0.01}
\end{align*}
Soft floor $\lfloor \cdot \rfloor_{0.01}$ prevents the habit term from
reaching zero (which would make the original equation singular).

\paragraph{Eq.~6: Bond Euler (Fisher equation).}
\begin{equation}
  -\textcolor{orange!80!black}{\lambda_{z,t}}
  + \textcolor{brown}{R_t}\,\beta\;\mathds{E}_t\!\left[
    \frac{\textcolor{blue}{\lambda_{z,t+1}}}
      {\textcolor{blue}{\pi_{t+1}}\,\textcolor{magenta}{\mu_{z,t+1}}}
  \right] = 0
\end{equation}

\paragraph{Eq.~7: Investment Euler (Tobin's Q).}
\begin{equation}
  \begin{split}
  0 = {}&
    \textcolor{orange!80!black}{\lambda_{z,t}}\,\textcolor{orange!80!black}{q_t}\!\left[
      1 - \textcolor{brown}{S(\cdot)}
      - \frac{\textcolor{green!60!black}{\mu_{z,t}}\,\textcolor{orange!80!black}{i_t}}
        {\textcolor{green!60!black}{i_{t-1}}}\;\textcolor{brown}{S'(\cdot)}
    \right]
    - \frac{\textcolor{orange!80!black}{\lambda_{z,t}}}{\textcolor{green!60!black}{\mu_{\Upsilon,t}}}
  \\
    &+ \beta\;\mathds{E}_t\!\left[
      \textcolor{blue}{\lambda_{z,t+1}}\,\textcolor{blue}{q_{t+1}}\,\textcolor{magenta}{\mu_{z,t+1}}
      \left(\frac{\textcolor{blue}{i_{t+1}}}{\textcolor{orange!80!black}{i_t}}\right)^{\!2}
      S'\!\!\left(\frac{\textcolor{magenta}{\mu_{z,t+1}}\,\textcolor{blue}{i_{t+1}}}
        {\textcolor{orange!80!black}{i_t}}\right)
    \right]
  \end{split}
\end{equation}

\paragraph{Eq.~8: Entrepreneur optimal contract.}
(Formerly eq.~9; eq.~8 bank participation is now eliminated analytically.)

Using next-period $\bar\omega_{t+1}$ from the Newton solver applied to next-period definitions:
\begin{equation}
  \begin{split}
  0 ={}&
    \frac{\textcolor{brown}{R_{t+1}^k}}{\textcolor{brown}{R_t}}
    \bigl(1 - \Gamma(\textcolor{brown}{\bar\omega_{t+1}})\bigr)
  \\
    &-\frac{\Gamma'(\textcolor{brown}{\bar\omega_{t+1}})}
      {\Gamma'(\textcolor{brown}{\bar\omega_{t+1}}) - \mu\,G'(\textcolor{brown}{\bar\omega_{t+1}})}
    \left[
      1 - \frac{\textcolor{brown}{R_{t+1}^k}}{\textcolor{brown}{R_t}}
      \bigl(\Gamma(\textcolor{brown}{\bar\omega_{t+1}}) - \mu\,G(\textcolor{brown}{\bar\omega_{t+1}})\bigr)
    \right]
  \end{split}
\end{equation}

\paragraph{Eq.~9: Resource constraint.}
$c_t$ is a network output; this equation enforces goods market clearing:
\begin{equation}
  0 = \textcolor{orange!80!black}{c_t} - \textcolor{brown}{c_t^{\text{implied}}}
\end{equation}

\medskip\noindent\textbf{Equation--output coupling.}\;
Every network output appears in at least one residual with nonzero gradient:
\begin{center}
\small
\begin{tabular}{ll}
\toprule
Output & Primary equations \\
\midrule
$\lambda_{z,t}$ & Eq.~1, 2b, 5, 6, 7 \\
$i_t$ & Eq.~7, 9 (via $k_t$, $y_z$, $c^{\text{impl}}$) \\
$\pi_t$ & Eq.~2a (via $K_{p,\text{inner}}$), Eq.~6 (via $R$), Eq.~9 \\
$\tilde{w}_t$ & Eq.~3, 4a, 9 (via $s$, $y_z$) \\
$h_t$ & Eq.~3, 4b, 9 (via $y_z$) \\
$F_{w,t}$ & Eq.~3, 4a \\
$F_{p,t}$ & Eq.~1, 2a \\
$q_t$ & Eq.~7, 9 (via $L$, $n$, monitoring) \\
$c_t$ & Eq.~5, 6 (via $y^{\text{GDP}} \to R$), 9 \\
$K_{p,t}$ & Eq.~2a, 2b \\
$K_{w,t}$ & Eq.~4a, 4b \\
\bottomrule
\end{tabular}
\end{center}
No output has zero gradient in the loss --- the system is fully determined.


%==========================================================================
\section{Analytical Eliminations}
\label{sec:eliminations}
%==========================================================================

Three equilibrium equations are static (no expectations) and can be solved
analytically, reducing the system from 12 controls and 12 equations to
\textbf{11 controls and 11 equations} (with $K_p$, $K_w$ as outputs and
the Phillips $K$ equations split).

\begin{center}
\begin{tabular}{clll}
\toprule
\# & Eliminated & From & Method \\
\midrule
1 & $s_t$ (marginal cost) & Eq.~10 (cost minimization FOC) & Direct solve \\
2 & $L_t$ (leverage) & Eq.~12 (balance sheet identity) & Direct solve \\
3 & $\bar\omega_t$ (default threshold) & Eq.~8 (bank participation) & Newton solver \\
\bottomrule
\end{tabular}
\end{center}

\noindent\textbf{Why $c_t$ is not eliminated.}\;
See Section~1.2 (policy variables).


%==========================================================================
\section{Laws of Motion}
%==========================================================================

\subsection{Exogenous states}
\begin{align*}
  \log\epsilon_t &= \rho_\epsilon \log\epsilon_{t-1} + \sigma_\epsilon\,e_t^\epsilon
  \\
  \log\mu_{\Upsilon,t} &= \rho_{\mu_\Upsilon} \log\mu_{\Upsilon,t-1}
    + \sigma_{\mu_\Upsilon}\,e_t^{\mu_\Upsilon}
  \\
  \log\mu_{z,t} &= \rho_{\mu_z} \log\mu_{z,t-1}
    + (1-\rho_{\mu_z})\log\mu_{z,\text{ss}}
    + \sigma_{\mu_z}\,e_t^{\mu_z}
  \\
  \log g_t &= \rho_g \log g_{t-1} + (1-\rho_g)\log g_{\text{ss}}
    + \sigma_g\,e_t^g
  \\
  m_t^p &= \sigma_p\,e_t^p
\end{align*}

\subsection{Endogenous state updates}

States are updated via soft clipping (softplus-based, gradient $\approx 1$ in
valid range, decays smoothly outside):
\begin{equation*}
  x_{\text{next}} = \ell + \frac{1}{s}\log(1+e^{s(x-\ell)})
  \quad\text{(lower clip)},\qquad
  x_{\text{next}} = u - \frac{1}{s}\log(1+e^{s(u-x)})
  \quad\text{(upper clip)}
\end{equation*}
with $s=5$. Applied to endogenous states only; exogenous states (AR(1)
processes) are not clipped.
\begin{equation*}
  \pi_{t-1} \leftarrow \pi_t, \quad
  k_{t-1} \leftarrow k_t, \quad
  c_{t-1} \leftarrow c_t, \quad
  q_{t-1} \leftarrow q_t, \quad
  i_{t-1} \leftarrow i_t, \quad
  R_{t-1} \leftarrow R_t, \quad
  \tilde{w}_{t-1} \leftarrow \tilde{w}_t, \quad
  L_{t-1} \leftarrow L_t
\end{equation*}


%==========================================================================
\section{Calibration}
\label{sec:calibration}
%==========================================================================

\begin{table}[h!]
\centering
\caption{Parameter calibration (quarterly frequency)}
\label{tab:calibration}
\begin{tabular}{lcll}
\toprule
Parameter & Symbol & Value & Source \\
\midrule
Discount rate & $\beta$ & 0.9985 & CMR calibration \\
Habit parameter & $b$ & 0.74 & CMR estimation \\
Inverse Frisch elasticity & $\sigma_L$ & 1.0 & CMR calibration \\
Disutility weight on labor & $\psi_L$ & 0.7705 & CMR calibration \\[3pt]
Capital share & $\alpha$ & 0.4 & CMR calibration \\
Depreciation rate & $\delta$ & 0.025 & CMR calibration \\
Adjustment cost & $\kappa$ & 2.0 & CMR: 10.78 \\
Fixed cost & $\Phi$ & 0.606 & CMR steady state \\[3pt]
Price markup & $\lambda_f$ & 1.2 & CMR calibration \\
Wage markup & $\lambda_w$ & 1.2 & CMR: 1.05 \\
Calvo price stickiness & $\xi_p$ & 0.6 & CMR: 0.74 \\
Calvo wage stickiness & $\xi_w$ & 0.6 & CMR: 0.81 \\
Price index weight & $\iota$ & 0.9 & CMR calibration \\
Wage index weight & $\iota_w$ & 0.49 & CMR calibration \\
Wage $\mu_z$ index weight & $\iota_\mu$ & 0.94 & CMR calibration \\[3pt]
Taylor rule smoothing & $\rho_p$ & 0.85 & CMR estimation \\
Taylor rule inflation & $\alpha_\pi$ & 1.5 & CMR: 2.4 \\
Taylor rule output & $\alpha_y$ & 0.36 & CMR estimation \\[3pt]
Tax on consumption & $\tau_c$ & 0.047 & CMR calibration \\
Tax on capital income & $\tau_k$ & 0.32 & CMR calibration \\
Tax on labor income & $\tau_l$ & 0.24 & CMR calibration \\[3pt]
Entrepreneur consumption & $\Theta$ & 0.005 & CMR steady state \\
Entrepreneur survival & $\gamma_e$ & 0.985 & CMR calibration \\
Entrepreneur transfer & $w^e$ & 0.005 & CMR calibration \\
Idiosyncratic dispersion & $\sigma_\omega$ & 0.26822 & CMR steady state \\
Monitoring cost & $\mu$ & 0.22 & CMR calibration \\[3pt]
SS inflation target & $\pi_{\text{ss}}$ & 1.006 & CMR calibration \\
SS growth rate & $\mu_{z,\text{ss}}$ & 1.0041 & CMR calibration \\
SS interest rate & $R_{\text{ss}}$ & 1.011678 & Steady state \\
SS output & $y_{\text{ss}}$ & 3.0308 & Steady state \\
SS government spending & $g_{\text{ss}}$ & 0.616 & CMR calibration \\[3pt]
$\rho_\epsilon$ / $\sigma_\epsilon$ & & 0.809 / 0.0046 & CMR estimation \\
$\rho_{\mu_\Upsilon}$ / $\sigma_{\mu_\Upsilon}$ & & 0.987 / 0.004 & CMR estimation \\
$\rho_{\mu_z}$ / $\sigma_{\mu_z}$ & & 0.146 / 0.00715 & CMR estimation \\
$\rho_g$ / $\sigma_g$ & & 0.94 / 0.023 & CMR estimation \\
$\sigma_p$ & & 0.0049 & CMR estimation \\
\bottomrule
\end{tabular}
\end{table}


%==========================================================================
\section{Steady State}
\label{sec:steady_state}
%==========================================================================

Numerically solved via \texttt{scipy.optimize.root} with JAX autodiff Jacobian.
All 11 residuals $< 6\times 10^{-8}$.

\begin{table}[h!]
\centering
\begin{minipage}[t]{0.45\linewidth}
\centering
\caption{Policy values at SS}
\begin{tabular}{lrl}
\toprule
Variable & Value & Bounding \\
\midrule
$\lambda_z$  & 0.6018 & softplus \\
$i$          & 0.7947 & softplus \\
$\pi$        & 1.0123 & sigmoid \\
$\tilde{w}$  & 1.9202 & softplus \\
$h$          & 0.9440 & softplus \\
$F_w$        & 0.8853 & softplus \\
$F_p$        & 4.7359 & softplus \\
$q$          & 1.0000 & softplus \\
$c$          & 1.5937 & softplus \\
$K_p$        & 4.8308 & softplus \\
$K_w$        & 2.2071 & softplus \\
\midrule
$\bar\omega$ & 0.4885 & Newton \\
$s$          & 0.8330 & analytical \\
$L$          & 1.9657 & analytical \\
\bottomrule
\end{tabular}
\end{minipage}\hfill
\begin{minipage}[t]{0.45\linewidth}
\centering
\caption{State values at SS}
\begin{tabular}{lr}
\toprule
State & Value \\
\midrule
$\pi_{t-1}$        & 1.0123 \\
$k_{t-1}$          & 27.421 \\
$c_{t-1}$          & 1.5937 \\
$q_{t-1}$          & 1.0000 \\
$i_{t-1}$          & 0.7947 \\
$R_{t-1}$          & 1.0180 \\
$\tilde{w}_{t-1}$  & 1.9202 \\
$L_{t-1}$          & 1.9657 \\
\midrule
$\epsilon_t$       & 1.0000 \\
$\mu_{\Upsilon,t}$ & 1.0000 \\
$g_t$              & 0.6160 \\
$\mu_{z,t}$        & 1.0041 \\
$m_t^p$            & 0.0000 \\
\bottomrule
\end{tabular}
\bigskip

\textbf{Note:} $\pi_{\text{SS}} = 1.012 \neq \pi_{\text{ss}} = 1.006$.
The Taylor rule with $\alpha_y > 0$ and financial frictions shifts
the equilibrium inflation above the target.
\end{minipage}
\end{table}


%==========================================================================
\section{Errata Summary}
\label{sec:errata}
%==========================================================================

The following bugs were found and corrected relative to the original
implementation. See \texttt{docs/disaster\_model\_errata.md} for details.

\begin{enumerate}
\item \textbf{Net worth formula} (Bug~1): used $(1-\mu G)$ instead of
  $(1-\Gamma)$, producing $n \approx 2\times$ too large.

\item \textbf{Tautological eq.~12} (Bug~2): capital accumulation identity
  was always zero (already enforced by dynamics). Replaced with balance
  sheet identity $L\,n = q\,k$.

\item \textbf{Inconsistent steady state} (Bug~3): hardcoded values had
  $O(10^{-2})$ residuals in wage Phillips equations. Replaced with
  numerical solver.
\end{enumerate}

Additionally, four errors in \texttt{disaster.tex} were caught during
the original code translation (the code has always been correct):

\begin{enumerate}
\setcounter{enumi}{3}
\item \textbf{Eq.~8 in tex}: uses default probability $F$ instead of
  survival probability $(1-F)$.

\item \textbf{Eq.~9 in tex}: uses an intermediate algebraic form; the
  code uses the numerically more stable rearrangement.

\item \textbf{Production function in tex}: writes leverage $L$ as labor
  input; the code correctly uses hours $h$.

\item \textbf{$G'(\bar\omega)$ in tex}: wrong sign ($+\sigma^2/2$
  instead of $-\sigma^2/2$) and missing $1/\bar\omega$ in denominator.
\end{enumerate}


\end{document}
